%----------------------------------------------------------------------------------------
% METADATA BLOCK (YAML)  — see inbox/README.md
%----------------------------------------------------------------------------------------
% uid: ""
% school: "lycee"
% level: "3eme"
% section: "Economie et Gestion"
% chapter: "Statistiques"
% title: "Trouver l'erreur"
% difficulty: "Facile"
% tags: ["statistiques", "recherche-erreur", "quartiles", "variance"]
% provenance: "own"
% batch: "2026-11-19-3eco-statistiques-n5"
%----------------------------------------------------------------------------------------

\newcommand{\R}{\mathbb{R}}
\newcommand{\vect}[1]{\overrightarrow{#1}}

% ---STATEMENT---
On donne la série de dix notes : \quad $8\quad12\quad7\quad15\quad10\quad9\quad14\quad11\quad10\quad14$.

Un élève a écrit :
\begin{itemize}
  \item Moyenne : $\overline x=\dfrac{110}{10}=11$.
  \item Série rangée : $7\ \ 8\ \ 9\ \ 10\ \ 10\ \ 11\ \ 12\ \ 14\ \ 14\ \ 15$.
  \item Médiane : la valeur de rang $5$, donc $\text{Me}=10$.
  \item Quartiles : $Q_1=\dfrac{x_2+x_3}{2}=8{,}5$ et $Q_3=\dfrac{x_7+x_8}{2}=13$.
  \item Variance : $V=\dfrac{7^2+8^2+\dots+15^2}{10}=127{,}6$.
\end{itemize}
\begin{enumerate}
  \item Quels résultats sont justes ? Repérer les trois erreurs.
  \item Corriger : donner la médiane, $Q_1$, $Q_3$, la variance et l'écart type (au centième).
\end{enumerate}

% ---SOLUTION---
\begin{enumerate}[itemsep=14pt]
  \item La moyenne et la série rangée sont justes : $7+8+9+10+10+11+12+14+14+15=110$.

  \textbf{Erreur 1 :} la médiane. Avec $n=10$ (pair), il faut les \emph{deux} valeurs centrales.

  \textbf{Erreur 2 :} les quartiles. Le manuel ne fait pas de moyenne de deux valeurs pour $Q_1$ et $Q_3$.

  \textbf{Erreur 3 :} la variance. L'élève a oublié de retrancher $\overline x^{\,2}$.

  \item
  \begin{prop}[Médiane et quartiles (convention du manuel)]
  $n$ pair : \ib{$\text{Me}=\dfrac{x_{n/2}+x_{n/2+1}}{2}$}. \ $Q_1$ : rang $\dfrac n4$ ou, s'il n'est pas entier, le rang immédiatement supérieur ; idem pour $Q_3$ avec $\dfrac{3n}4$.
  \end{prop}
  \[\text{Me}=\dfrac{x_5+x_6}{2}=\dfrac{10+11}{2}=10{,}5\]
  $\dfrac{10}{4}=2{,}5$ : $Q_1$ est de rang $3$, donc $Q_1=9$.

  $\dfrac{30}{4}=7{,}5$ : $Q_3$ est de rang $8$, donc $Q_3=14$.
  \begin{prop}[Variance]
  \ib{$V=\text{(moyenne des carrés)}-\overline x^{\,2}$}.
  \end{prop}
  \[V=\dfrac{1\,276}{10}-11^2=127{,}6-121=6{,}6\qquad\sigma=\sqrt{6{,}6}\approx2{,}57\]
  \textbf{Contrôle :} une variance de $127{,}6$ donnerait $\sigma\approx11{,}3$, soit autant que la moyenne : impossible pour des notes entre $7$ et $15$.
\end{enumerate}
